\begin{afterword}
\summaryhead

The post defines anthropomorphism as expecting a process unlike a human brain to do what your brain predicts. Its example is the group selectionists, biologists before 1966 who believed predators would restrain their breeding for the good of the group. When Michael Wade selected insect populations for small size in the laboratory, the insects instead cannibalized eggs and larvae.

The author explains why the biologists missed this. A brain searches for solutions that rank high in its own preferences, and does not even generate repugnant ones like cannibalism. Evolution ranks outcomes by a different standard. The author argues that the biologists kept the conclusion that evolution would do what they would do, and replaced the open reasoning with rationalizations, an instinct humans may have evolved for tribal politics. The post notes that this has a subtext about artificial intelligence, and generalizes: optimism is ranking outcomes by your own preferences and treating the best as a prediction.

\responsehead

The general warning is sound, and the post states it in a modest form: a process that works differently from you will ``not necessarily'' do what you would do. The problems are the example, on which almost all the argument depends, and the account of the biologists' minds, which is asserted rather than shown.

The post claims to know how the group selectionists reasoned. It claims that they ``erased what had been written above the bottom line'' and ``wrote in new rationalizations.'' No work by Wynne-Edwards, Allee or Brereton is quoted. The process is inferred from the conclusion: the conclusion was mistaken and pleasing, so the reasoning must have been rationalization. The post's last paragraph tells the reader to ``look at the cognitive history''; for its own example, it does not.

The published record does not fit the picture well, as the notes on ``The Tragedy of Group Selectionism'' show: Wynne-Edwards described competition and death, not ``quiet peace,'' and cannibalism was a known check on flour-beetle numbers.

In short: a sound warning against predicting other processes by your own preferences, illustrated by a history of the group selectionists' reasoning that quotes none of them and fits the written record poorly.
\end{afterword}
